\section{三次样条：用局部性换稳定}
\label{sec:08-Numerical-Analysis/03-Interpolation-and-Approximation/04}

数据录入的时候你把一个数打错了。发现之后改回来，重新拟合——曲线在哪些地方会变？

用全局多项式，答案是``到处''。那个错值参与了每一个系数的计算，改动它，整条曲线从最左端到最右端一起挪窝。用样条，答案是``错点附近那几段''，往外几个节点影响就衰减到看不见了。

这个差别不是实现细节，它决定了模型可不可以被信任地维护。数据集会打补丁、会有异常值被剔除、会有新样本追加。牵一发动全身的模型每次都要整体复检；影响可局部化的模型只需检查邻域。

\begin{center}
\begin{tikzpicture}[x=0.95cm,y=0.92cm,font=\small,>=stealth]
\begin{scope}
  \draw[->,black!55] (-0.2,0) -- (6.5,0) node[right,font=\footnotesize] {\(x\)};
  \draw[->,black!55] (0,-0.2) -- (0,3.5);
  \draw[black!45,thick] plot[smooth] coordinates {
    (0.0,0.400) (0.4,1.094) (0.8,1.544) (1.2,1.520) (1.6,1.187) (2.0,1.100)
    (2.4,1.618) (2.8,2.255) (3.2,2.342) (3.6,1.871) (4.0,1.500) (4.4,1.703)
    (4.8,2.106) (5.2,2.143) (5.6,1.657) (6.0,0.900)};
  \draw[blue!75,thick] plot[smooth] coordinates {
    (0.0,0.400) (0.4,1.056) (0.8,1.510) (1.2,1.576) (1.6,1.326) (2.0,1.100)
    (2.4,1.150) (2.8,1.333) (3.2,1.420) (3.6,1.403) (4.0,1.500) (4.4,1.841)
    (4.8,2.161) (5.2,2.110) (5.6,1.618) (6.0,0.900)};
  \foreach \p in {(0,0.4),(1,1.6),(2,1.1),(4,1.5),(5,2.2),(6,0.9)}
    \fill[black!80] \p circle (1.6pt);
  \fill[black!80] (3,2.4) circle (1.6pt);
  \fill[blue!80] (3,1.4) circle (1.6pt);
  \draw[blue!60,->] (3,2.28) -- (3,1.52);
  \node[font=\footnotesize,align=center] at (3.0,-0.85) {三次样条\\影响被关在邻近几段};
\end{scope}
\begin{scope}[shift={(8.2,0)}]
  \draw[->,black!55] (-0.2,0) -- (6.5,0) node[right,font=\footnotesize] {\(x\)};
  \draw[->,black!55] (0,-0.2) -- (0,3.5);
  \draw[black!45,thick] plot[smooth] coordinates {
    (0.0,0.400) (0.4,2.779) (0.8,2.153) (1.2,1.137) (1.6,0.762) (2.0,1.100)
    (2.4,1.767) (2.8,2.294) (3.2,2.381) (3.6,2.021) (4.0,1.500) (4.4,1.275)
    (4.8,1.720) (5.2,2.757) (5.5,3.355) (5.7,3.173) (6.0,0.900)};
  \draw[blue!75,thick] plot[smooth] coordinates {
    (0.0,0.400) (0.4,1.789) (0.8,1.780) (1.2,1.409) (1.6,1.145) (2.0,1.100)
    (2.4,1.208) (2.8,1.347) (3.2,1.434) (3.6,1.461) (4.0,1.500) (4.4,1.658)
    (4.8,1.992) (5.2,2.384) (5.5,2.453) (5.7,2.190) (6.0,0.900)};
  \foreach \p in {(0,0.4),(1,1.6),(2,1.1),(4,1.5),(5,2.2),(6,0.9)}
    \fill[black!80] \p circle (1.6pt);
  \fill[black!80] (3,2.4) circle (1.6pt);
  \fill[blue!80] (3,1.4) circle (1.6pt);
  \draw[blue!60,->] (3,2.28) -- (3,1.52);
  \node[font=\footnotesize,align=center] at (3.0,-0.85) {全局六次插值\\整条曲线重排};
\end{scope}
\end{tikzpicture}

\emph{同样七个节点，把 \(x=3\) 处的数据下移一格（箭头），灰线是改动前、蓝线是改动后。左边的样条在 \(x=0.5\) 处只动了 \(0.04\)，占改动量的 \(4\%\)；右边的全局多项式在同一位置动了 \(0.90\)，占 \(90\%\)，而且区间另一端也跟着大幅改写。}
\end{center}

\subsection{把区间切开，每段只用三次}

样条的定义就是把这个想法写清楚。给定严格递增的节点 \(x_0<\cdots<x_n\)，三次插值样条 \(s\) 在每个 \([x_i,x_{i+1}]\) 上是次数不超过 \(3\) 的多项式，满足 \(s(x_i)=y_i\)，并且 \(s\)、\(s'\)、\(s''\) 在整个区间上连续。

为什么是三次？次数再低就接不平：分段线性只有 \(C^0\)，折角处曲率无定义；分段二次能做到 \(C^1\)，但二阶导仍会跳，曲线看上去会有``忽然转向''的观感。三次刚好让每段的四个自由度够用——两端函数值加上两端的一阶、二阶导衔接，同时保证方程组仍是带状的。次数再往上加，自由度过剩，局部性也会被稀释。

数一数自由度就知道还差什么。\(n\) 段三次式共 \(4n\) 个系数；插值条件（每段两端各一个）用掉 \(2n\) 个，\(n-1\) 个内节点上的一阶、二阶导连续再用掉 \(2(n-1)\) 个，还剩 \(2\) 个。这两个自由度必须由端点条件补上，常见的做法是令两端二阶导为零（natural）、指定两端斜率（clamped），或者要求首尾两个内节点不真正断开（not-a-knot）。

这三种不是可以随手互换的语法糖，它们代表你对端点掌握了什么信息。知道物理上的端点斜率就该用 clamped；natural 条件并不意味着真实函数在端点曲率为零，它只是在你什么都不知道时挑了一个说法。选错的后果集中在端点附近，端点误差往往会拖低整体的收敛阶。所以报告结果时``用了三次样条''这句话是不完整的，必须连边界条件一起说。

\subsection{方程组为什么是三对角的}

推导可以只留骨架。以节点二阶导 \(M_i=s''(x_i)\) 为未知量，记 \(h_i=x_{i+1}-x_i\)。在每段上 \(s''\) 是线性的（三次式的二阶导），由两端值 \(M_i,M_{i+1}\) 完全确定；积两次分、用两端函数值定住积分常数，该段的表达式就写出来了，两端的一阶导也随之可算。再要求一阶导在每个内节点左右相等，整理即得

\[
h_{i-1}M_{i-1}+2(h_{i-1}+h_i)M_i+h_iM_{i+1}
=6\left(\frac{y_{i+1}-y_i}{h_i}-\frac{y_i-y_{i-1}}{h_{i-1}}\right).
\]

关键不在这个式子的样子，而在它只牵涉相邻三个未知量。每个方程一条，凑成的矩阵除主对角线和上下两条副对角线外全是零。这种三对角系统在节点严格递增、边界条件合法时严格对角占优，可解且解唯一，用追赶法 \(O(n)\) 次运算解完——对比一般稠密系统的 \(O(n^3)\)，差距随规模拉开。求值同样便宜：先定位查询点落在哪一段，再算一个三次式。

带状结构正是局部性在代数层面的样子。改一个 \(y_i\)，右端只有三个分量变化；扰动沿三对角系统传播时按几何速率衰减，隔开几个节点就已经不可见。上面那张图的 \(4\%\) 对 \(90\%\)，就是这条性质的数值版本。

\subsection{``看起来平滑''可以被算出来}

样条给人的第一印象是曲线好看，这个印象有个精确的对应物。

\begin{theorem}
在所有通过给定数据点、二阶弱可导的函数 \(g\) 中，natural 三次样条 \(s\) 使弯曲能量 \(\int_{x_0}^{x_n}[g''(x)]^2dx\) 取到最小。
\end{theorem}

证明骨架短得可以写在这里：令 \(g=s+h\)，则 \(h\) 在所有节点为零。展开平方后的交叉项 \(\int s''h''\) 逐段分部积分，段内那部分因 \(s'''\) 在段内是常数、\(h\) 在节点处为零而消失，端点那部分因 natural 条件 \(s''=0\) 而消失。于是 \(\int(g'')^2=\int(s'')^2+\int(h'')^2\ge\int(s'')^2\)。条件对账也很直接：natural 边界是端点项归零的前提，换成 clamped 会得到另一个版本的最小性陈述，而不是同一条。

顺带说一句名字的来历。造船和绘图用的样条是一根有弹性的细木条，用压铁固定在若干点上，木条自然弯成的形状使弹性势能最小——在小挠度下，弹性势能正比于 \(\int(g'')^2\)。数学定义和那根木条对得上。

\subsection{局部支撑：这笔交易后面还会做很多次}

三次样条是本书里第一次出现``局部支撑''这个念头：每个基函数只在有限个区间上非零，因此每个系数只对局部负责。把它写成基的语言就是 B 样条——一族只覆盖相邻几个区间的分段多项式，样条空间中的任何函数都是它们的组合。

同一笔交易此后反复出现。有限元把区域剖成单元，每个基函数只支撑在一个节点周围的几个单元上，于是刚度矩阵是稀疏的。小波用有限支撑的母函数做多尺度分解，代价是频率分辨率，换来的是时间定位能力，见\hyperref[sec:08-Numerical-Analysis/08-Fourier-and-Spectral-Methods/06]{``小波与多尺度分解''}。卷积核只看感受野内的像素，注意力的滑动窗口只连接邻近的 token——它们放弃了全局连接的表达力，换来了参数量可控、扰动可局部化、以及``这个输出受哪些输入影响''这个问题有明确答案。

代价是要说清楚的。局部基不擅长表达真正的全局结构：一个在整个区间上缓慢起伏的趋势，用局部基去拼要花很多段，而一个恰当的全局基一两项就够。全局与局部之间没有普适的最优，只有对当前信号的判断——你要刻画的东西是分布在整体上还是集中在局部。

\subsection{光滑不等于保形，也不等于能外推}

三次样条会做几件让人意外的事。数据单调上升时，样条可能在两点之间先冲高再回落；数据全为正时，样条可能探到负值。\(C^2\) 管的是导数连续，它不承诺单调、不承诺凸性、不承诺正性。画一条浓度曲线或概率曲线时，负值是会被用户当成错误报上来的。

要保形就得换方法。PCHIP 一类的分段三次 Hermite 插值牺牲二阶光滑，只保 \(C^1\)，换来单调性的保持。选之前先想清楚需要的是哪条性质，而不是默认越光滑越好。

收敛阶也有前提。函数足够光滑、网格受控、边界条件与真实端点行为相容时，三次样条的函数值误差可达 \(O(h^4)\)。这三个条件缺一不可：natural 条件配上端点曲率非零的真函数，端点附近的精度会明显掉一档；网格严重不均或函数不够光滑，常数和阶都可能改变。做网格加密实验时若看到的阶不是 \(4\)，先检查这几条，再怀疑代码。

最后是外推。走出最后一个节点，曲线的走向完全由那一段的三次式和边界条件决定，与数据的整体趋势无关。natural 条件下末段的二阶导为零，外推近似是一条直线；换个边界条件就是另一个故事。样条在数据范围内的稳定性，一步都不外借。

数据含噪时还有一条岔路：精确插值样条仍然会逐点追踪噪声，此时该用的是平滑样条——在弯曲能量和残差之间加一个权衡参数，让曲线在``贴近数据''和``别太扭''之间取舍。那个参数扮演的角色，和岭回归里的正则系数是同一个，见\hyperref[sec:13-Machine-Learning/02-Linear-Models/02]{``正则化从 Ridge 到 Lasso''}。

到这里，插值与逼近的工具箱基本齐了：怎样穿过数据、什么时候不该穿过、节点放在哪里、区间该不该切开。下一章把这些曲线拿去做两件相反的事——求导和求积分。同一份采样，积分越加越准，求导却会把噪声放大，理由就藏在本章反复出现的那个误差结构里，见\hyperref[ch:064]{``数值微分与积分：同一份采样，为何难度相反''}。
